\section{Objetivo}

Este trabalho documenta, com evidências de execução, o sistema de supervisão \textit{Lean} desenvolvido para uma linha didática de 12 postos. O sistema integra um CLP LS XGK-CPUE, que executa o programa Ladder da linha; uma Raspberry Pi, que atua como \textit{gateway} de borda entre a porta serial do CLP e a rede IP e comanda o motor da esteira; um receptor no Windows, que persiste os eventos e serve a Dashboard web; a Instrução de Trabalho Digital (IT Digital) do posto E07; e um aplicativo Desktop que simula os 12 postos quando o CLP não está disponível. Os objetivos são: (i) descrever a arquitetura e os protocolos de cada enlace; (ii) documentar cada componente - Desktop, Dashboard, Raspberry Pi, CLP, IT Digital e APIs - a partir do código-fonte e da configuração efetivamente usados; e (iii) apresentar, para cada componente, evidências concretas de funcionamento: registros de execução com o CLP real, bancos de dados, testes automatizados, chamadas reais às APIs e capturas de tela.

\subsection{Materiais e reprodução}

O código está organizado em duas aplicações Python. \textbf{Gateway/} contém o serviço da Raspberry Pi (Modbus RTU, OPC UA, TCP, SNTP, SQLite e GPIO). \textbf{MockDesktop/} contém o simulador Desktop (PyQt6), o receptor TCP de referência, o cliente OPC UA e a Dashboard web (FastAPI). A documentação do programa Ladder está em \textbf{Docs/DOCUMENTACAO\_CLP\_LS\_XGK.md}, com 21 capturas do projeto \cod{LAB\_NEW\_UPDATED} no XG5000 4.74. Ambiente utilizado nas verificações: Windows 11, Python 3.12.10, PyQt6 6.8.1, asyncua 1.1.5, pymodbus 3.11.3, FastAPI 0.115.11, uvicorn 0.34.0, cryptography 44.0.2 e pytest 8.3.5. Na Raspberry Pi o serviço usa Python 3.11 ou superior, gpiozero 2.0.1 e pigpio 1.78.

As evidências têm três origens, sempre identificadas no texto:

\begin{itemize}
    \item \textbf{Execução real com o CLP} - arquivo \cod{Gateway/logs/gateway.log} (32.140 linhas, de 27/09/2026 16:57 a 28/09/2026 13:11) e banco \cod{Gateway/data/gateway.db} (18 eventos de Andon capturados no CLP pelo Gateway \cod{LINHA\_A}), além do banco do receptor \cod{MockDesktop/data/reference\_receiver.db} (127 recepções entre 27/09 e 21/11/2026).
    \item \textbf{Testes automatizados} - 98 testes do MockDesktop e 66 do Gateway, executados em 21/11/2026 (seção 9).
    \item \textbf{Demonstração ponta a ponta com CLP emulado} - para gerar capturas de tela e medições sem o hardware, o Gateway, o receptor e a Dashboard reais foram executados contra um servidor Modbus RTU que reproduz o mapa de registradores do Ladder (\cod{clp\_emulado.py}, descrito na seção 4.6). Os resultados dessa demonstração estão sempre rotulados como ``CLP emulado''. Os bancos e portas da demonstração foram separados dos da instalação (portas 4850, 8091 e 15010), de modo que os dados reais não foram alterados.
\end{itemize}
